\documentclass[12pt]{article}

\usepackage[a4paper, margin=1in]{geometry}
\usepackage{amsmath,amssymb,amsthm}
\usepackage{mathtools}
\usepackage{mathrsfs}
\usepackage{array}
\usepackage{booktabs}
\usepackage{longtable}
\usepackage{tabularx}
\usepackage{enumitem}
\usepackage{float}
\usepackage{xcolor}
\usepackage[hidelinks]{hyperref}
\usepackage{fontspec}
\usepackage{xeCJK}

\defaultfontfeatures{Ligatures=TeX}
\IfFontExistsTF{Times New Roman}
  {\setmainfont{Times New Roman}}
  {\setmainfont{TeX Gyre Termes}}
\IfFontExistsTF{Menlo}
  {\setmonofont{Menlo}}
  {\setmonofont{Latin Modern Mono}}
\IfFontExistsTF{PingFang SC}
  {\setCJKmonofont{PingFang SC}}
  {\IfFontExistsTF{Songti SC}
    {\setCJKmonofont{Songti SC}}
    {}}
\IfFontExistsTF{PingFang SC}
  {\setCJKmainfont{PingFang SC}}
  {\IfFontExistsTF{Songti SC}
    {\setCJKmainfont{Songti SC}}
    {\setCJKmainfont{STSong}}}

\newtheorem{theorem}{Theorem}[section]
\newtheorem{proposition}[theorem]{Proposition}
\newtheorem{lemma}[theorem]{Lemma}
\newtheorem{corollary}[theorem]{Corollary}
\newtheorem{conjecture}[theorem]{Conjecture}
\theoremstyle{definition}
\newtheorem{definition}[theorem]{Definition}
\newtheorem{assumption}[theorem]{Assumption}
\newtheorem{example}[theorem]{Example}
\theoremstyle{remark}
\newtheorem{remark}[theorem]{Remark}

\newcolumntype{P}[1]{>{\raggedright\arraybackslash}p{#1}}
\newcolumntype{Y}{X}

\newcommand{\tightlist}{%
  \setlength{\itemsep}{0pt}\setlength{\parskip}{0pt}}

\newcounter{none}

\setlength{\LTleft}{0pt}
\setlength{\LTright}{0pt}


\title{%
  \textbf{Geometry Relative to a Cut:}\\
  \textbf{Selection, Generation and Modular Inseparability}%
}

\author{
  Sidong Liu, PhD \\
  \small iBioStratix Ltd \\
  \small \texttt{sidongliu@hotmail.com}
}

\date{October 2026}

\hypersetup{pdftitle={Geometry Relative to a Cut: Selection, Generation and Modular Inseparability},pdfauthor={Sidong Liu}}
\begin{document}
\emergencystretch=2em
\raggedbottom
\hbadness=5000
\vbadness=5000

\maketitle

\begin{abstract}
A perspective has a physical position in a process: its boundary determines which distinctions remain available, while its retained history can alter how those distinctions are used. We study the structural part of this relation through standard half-sided modular inclusions. A standard Cut and vacuum determine a strongly additive chiral net, and the net determines the modular triple. For a class of local stress seeds, the past modular orbit constructs a proper standard Cut without an independent forward-translation permission. In the SU(2) level-one model, an explicit random-rotation noise selects the same Virasoro algebra generated by the seed's ideal closure, with the two constructions taken in their appropriate vacuum representations.

The correspondence does not make geometric isomorphism a complete description of access. For a fixed local excitation, moving a Cut loses relative entropy by \(2\pi aE\). We show that this loss prevents a common physical recovery of the vacuum and excited state: the boundary change cannot be replaced by reversible recoding, even when the geometric triples remain isomorphic. Finite current-control programs give a task-level interface to the ideal stress flows.

We then connect this distinction to causal self-description. Matching content and confidence need not match the retained distinctions needed for correction; matching all permitted interventional record laws leaves no further test in that class. These results support a positive account of perspective as maintained, history-dependent access and specify selective interventions for its stronger reflective interpretation. That interpretation requires empirical evidence beyond the modular correspondence.
\end{abstract}

\section{Introduction}\label{sec:intro}

A view of the world depends on where and how a process can interact with it.
A detector's position matters, but so do its permitted operations, retained
records and capacity to change its own coupling. Two descriptions can use
different coordinates for the same experiment. Two physically different
access boundaries can instead make different distinctions available, even
inside the same geometric world. A theory of perspective must separate
these cases.

The question of experience remains broader. Chalmers distinguishes the
explanation of discrimination, internal access and report from the question
of why those functions are accompanied by experience \cite{Chalmers}.
The present paper investigates an organizing relation that such an
explanation has to locate: how the accessible world and the process's
conditions of access constrain one another. We call their treatment as
freely variable formal data the \emph{independence assumption}. This is a
specific mathematical question, not a premise attributed to every
formulation of the explanatory gap.

\paragraph{Four objects.}
An exact geometric boundary is a proper half-sided modular inclusion
\(N\subset M\), called a Cut. A noise-selected algebra is a persistent
subalgebra of a supplied carrier. Finite access is a menu of preparations,
instruments and controls at stated budgets. Reflective organization is a
further causal role: acquired information about the system's own access
can govern its later inquiry and correction. The first three can be
defined without attributing experience. Connecting the fourth to
first-person experience is the bridge hypothesis considered in
Section~\ref{sec:hard}, together with its competing explanations.

\paragraph{The geometric setting.}
Let \(\Omega\) be cyclic and separating for a von Neumann algebra \(M\).
Its modular group \(\sigma_s\) is fixed by the pair \((M,\Omega)\).
For a conformal half-line in its vacuum it implements dilations
\cite{BW,BGL}. A standard Cut reconstructs translations and a strongly
additive local conformal net, and conversely that net determines its
standard modular triple \cite{GLW}. This is a correspondence between
complete data, not a claim that a bare modular flow is a perspective.
It provides an exact structural setting in which geometry and a boundary
of access cannot be specified arbitrarily.

\paragraph{Selection, formation and response.}
Sections~\ref{sec:cut}--\ref{sec:response} establish three links.
First, a standard Cut carries the modular data of its geometric net.
Second, a specified random-rotation noise selects the Virasoro subnet
inside the SU(2) level-one carrier, and a bounded stress seed generates
that same algebra under ideal closure. For positive endpoint-monotone
logarithmic bumps, the seed's past modular orbit itself generates a
proper standard Cut (Theorem~\ref{thm:singleorbit}); its full orbit
generates the half-line (Corollary~\ref{cor:modonly}). For arbitrary
nonzero stress seeds, the dilation-plus-forward-translation result
remains available (Theorem~\ref{thm:hull}). Noise selection and the Cut
construction act at different representation levels, as
Remark~\ref{rem:twolevel} explains.

Third, a local excitation loses exactly \(2\pi aE\) of relative entropy
when the accessible half-line is displaced
(Theorem~\ref{thm:displacement}). Proposition~\ref{prop:access-recoding}
uses this loss to distinguish a physical restriction from a reversible
change of coding. The same isomorphism class of geometric triples can
therefore coexist with different access to a fixed excitation.
The location of a boundary relative to the source matters.

\paragraph{From a boundary to an organized perspective.}
The carrier, vacuum and modular dynamics are inputs to the construction;
it supplies a Cut for the specified seed class rather than deriving an
ambient net from a finite detector. Section~\ref{sec:limits} connects
constructed stress flows to finite current-control programs. Foundation~I
separates that chiral construction from finite central records, whose
stable hull need not grow at all \cite{FoundationI}.

Section~\ref{sec:hard} then asks what turns usable access into reflective
organization. Consciousness~III and \emph{The World as Imprint} supply
the relevant causal distinction: a record matters to this organization
when it changes the process's later permissions or corrections, rather
than merely describing them \cite{PaperCIII,WorldImprint}. Their results
also fix the limit of a discriminating test. Equal routine reports can
hide different causal histories; equal complete laws under all allowed
interventions cannot be separated by another report in the same class.
This permits a positive, testable bridge without identifying the exact
Cut, modular time or calibration success with experience.

\section{Preliminaries}\label{sec:prelim}

\subsection{Modular theory and half-sided modular inclusions}

Let \(M\) be a von Neumann algebra on a Hilbert space \(H\) and \(\Omega\in H\)
a vector that is cyclic and separating for \(M\). Write \(\Delta_M,J_M\) for the
modular operator and conjugation of \((M,\Omega)\) and
\(\sigma^M_s(x)=\Delta_M^{is}x\Delta_M^{-is}\). Then \(\sigma^M_s(M)=M\) and
\(\Delta_M^{is}\Omega=\Omega\).

\begin{definition}[Half-sided modular inclusion]\label{def:hsm}
Let \(N\subset M\) be von Neumann algebras for which \(\Omega\) is cyclic and
separating. The inclusion is \emph{\(+\)hsm} if \(\sigma^M_s(N)\subset N\) for all
\(s\le0\) (equivalently \(\Delta_M^{-it}N\Delta_M^{it}\subset N\) for \(t\ge0\)). It
is \emph{standard} if \(\Omega\) is also cyclic for \(M\cap N'\). It is \emph{proper}
if \(N\ne M\). We call a proper \(+\)hsm inclusion a \emph{Cut}.
\end{definition}

\begin{theorem}[Wiesbrock; Araki--Zsid\'o {\cite[Thm.~2.1]{AZ}}]\label{thm:AZ}
Let \(N\subset M\) be \(+\)hsm. Then \(P:=\overline{(\log\Delta_N-\log\Delta_M)/2\pi}\)
is a positive self-adjoint operator, and with \(U(s)=e^{isP}\):
\begin{enumerate}[label=(\arabic*),leftmargin=2em,itemsep=1pt]
\item \(\Delta_M^{-it}U(s)\Delta_M^{it}=U(e^{2\pi t}s)=\Delta_N^{-it}U(s)\Delta_N^{it}\);
\item \(J_MU(s)J_M=U(-s)=J_NU(s)J_N\);
\item \(\Delta_N^{it}=U(1)\Delta_M^{it}U(1)^*\) and \(J_N=U(1)J_MU(1)^*\);
\item \(N=U(1)MU(1)^*\);
\item \(U(s)MU(s)^*\subset M\) for \(s\ge0\), and \(\mathrm{Ad}\,U(s)\), \(s\ge0\), is
a strongly continuous semigroup of endomorphisms of \(M\);
\item \(U(s)\Omega=\Omega\), so the vacuum state is invariant under \(\mathrm{Ad}\,U(s)\).
\end{enumerate}
\end{theorem}

\begin{theorem}[{\cite[Thm.~2.2(4)]{AZ}}, {\cite[Cor.~1.8]{GLW}}]\label{thm:type}
If \(N\subset M\) is \(+\)hsm and \(p\) is the largest central projection of \(M\)
with \(Mp=Np\), then \(M(1-p)\) and \(N(1-p)\) are of type III\(_1\) whenever
\(p\ne1\). In particular, if \(M\) is a factor and \(N\ne M\), both are type III\(_1\)
factors.
\end{theorem}

\begin{theorem}[Guido--Longo--Wiesbrock {\cite[Cor.~1.7--1.9]{GLW}}]\label{thm:GLW}
There is a one-to-one correspondence between isomorphism classes of
\emph{standard} \(+\)hsm inclusions \((N\subset M,\Omega)\) and isomorphism
classes of strongly additive local conformal nets on \(S^1\); under it
\(M=\mathcal A(0,\infty)\), \(N=\mathcal A(1,\infty)\), \(\Omega\) is the vacuum,
\(U\) is the translation group, and \(\Delta_M^{it}\) acts as dilation. For a
conformal net that is not strongly additive, the hsm reconstruction returns the
dual net \(\mathcal A^d(a,b)=\mathcal A(a,\infty)\cap\mathcal A(b,\infty)'\), which
equals \(\mathcal A\) if and only if \(\mathcal A\) is strongly additive.
\end{theorem}

\subsection{Virasoro and current nets}

A conformal net \(\mathcal A\) is \emph{diffeomorphism covariant} with central
charge \(c>0\) if the vacuum representation carries a positive-energy projective
representation \(V\) of \(\mathrm{Diff}_+(S^1)\), with \(V(\varphi)\in\mathcal A(I)\)
when \(\varphi\) is supported in \(I\), and stress tensor \(T\) with
\(V(\mathrm{Exp}\,f)=e^{iT(f)}\) up to a phase (Goodman--Wallach \cite{GW}, Carpi--Weiner \cite{CW}).
Write \(\mathrm{Vir}_c(I)=\{T(f):\mathrm{supp}f\subset I\}''\). For real \(f\in
C_c^\infty\), \(T(f)\) is essentially self-adjoint on finite-energy vectors, and if
\(f_n\to f\) in \(\|f\|_{3/2}=\sum_n|\hat f_n|(1+|n|^{3/2})\) with supports in a fixed
compact set, then \(T(f_n)\to T(f)\) in the strong resolvent sense
\cite[Thm.~4.4, Prop.~4.5]{CW}. For the half-line \(M=\mathrm{Vir}_c(0,\infty)\) the
modular group of the vacuum is the dilation group, \(\sigma_s=\mathrm{Ad}\,V(\delta_{e^{-2\pi s}})\)
with \(\delta_\lambda(x)=\lambda x\) (Bisognano--Wichmann \cite{BW}; Brunetti--Guido--Longo \cite{BGL}), and \(T_a=\mathrm{Ad}\,U(a)\) with \(U(a)=V(\tau_a)\) implements
translations. The transformation law is
\(V(\varphi)T(x)V(\varphi)^*=\varphi'(x)^2T(\varphi(x))-\tfrac c{24\pi}\mathrm{Sch}(\varphi)(x)\).
Virasoro nets with \(c\le1\) are strongly additive and for \(c>1\) they are not
\cite[\S2.4]{Carpi}; the fixed-point subnet of a strongly additive net under a
compact group is strongly additive \cite{Xu}.

\subsection{Relative entropy of coherent states}\label{sec:hollands}

Let \(f\in C_c^\infty((a_0,\infty),\mathbb R)\) with \(a_0>0\), \(\rho=\mathrm{Exp}\,f\),
\(U=V(\rho)=e^{iT(f)}\), \(\omega_+(A)=\omega(U^*AU)\) and
\(\omega_-(A)=\omega(UAU^*)\). Put \(w=\rho''/\rho'\). Hollands
\cite[Thm.~1]{Hollands} proves, for the half-line algebra \(\mathrm{Vir}_c(0,\infty)\),
\begin{equation}\label{eq:hollands}
S(\omega\|\omega_+)=\frac{c}{24}\int_0^\infty x\,(\partial_x\log\rho'(x))^2\,dx .
\end{equation}

\section{Cuts and geometry}\label{sec:cut}

Let \((N\subset M,\Omega)\) be \(+\)hsm with translation group \(U\) from
Theorem~\ref{thm:AZ}. Put \(N_a=U(a)MU(a)^*\) for \(a\ge0\), so \(N_0=M\), \(N_1=N\).

\begin{lemma}[The modular orbit of a Cut generates \(M\)]\label{lem:orbit}
\(\bigvee_{t\in\mathbb R}\sigma^M_t(N)=M\).
\end{lemma}
\begin{proof}
By Theorem~\ref{thm:AZ}(1),(4) and \(\sigma_t(M)=M\),
\(\sigma_t(N)=\Delta_M^{it}U(1)MU(1)^*\Delta_M^{-it}=U(e^{-2\pi t})MU(e^{-2\pi t})^*
=N_{e^{-2\pi t}}\). For \(m\in M\) put \(m_a=U(a)mU(a)^*\in N_a\). As
\(t\to\infty\), \(a=e^{-2\pi t}\to0\) and \(m_a\to m\) strongly, since \(U\) is
strongly continuous. Each \(m_a\) lies in \(\sigma_t(N)\) for \(a=e^{-2\pi t}\), and
the algebra generated by the \(\sigma_t(N)\) is strongly closed, so it contains
\(m\). The other inclusion is trivial.
\end{proof}

\begin{lemma}[One-directional strict nesting]\label{lem:nest}
If \(N\subsetneq M\) is \(+\)hsm, then \(\sigma_t(N)\subsetneq N\) for \(t<0\) and
\(\sigma_t(N)\not\subset N\) for \(t>0\).
\end{lemma}
\begin{proof}
For \(a\ge b\ge0\), \(N_a=U(b)[U(a-b)MU(a-b)^*]U(b)^*\subset U(b)MU(b)^*=N_b\) by
Theorem~\ref{thm:AZ}(5), so \(a\mapsto N_a\) is decreasing. Suppose \(N_a=N_b\) for
some \(a>b\ge0\). Conjugating by \(U(b)^*\) gives \(N_d=M\) with \(d=a-b>0\). Choose
\(t\) with \(e^{-2\pi t}d=1\): then \(N=N_1=\sigma_t(N_d)=\sigma_t(M)=M\), a
contradiction. Hence \(a\mapsto N_a\) is strictly decreasing. Since
\(\sigma_t(N)=N_{e^{-2\pi t}}\), the claims follow.
\end{proof}

The nesting is a property of the subalgebras. It is not entropy production or an arrow
of time: the modular group itself is invertible and acts by automorphisms.
Lemma~\ref{lem:nest} says that a proper Cut cannot be stable in both directions.

\begin{proposition}[A Cut lives in type III\(_1\)]\label{prop:typeIII}
If \(M\) is a factor and \(N\subsetneq M\) is \(+\)hsm, then \(M\) and \(N\) are
type III\(_1\) factors. No proper \(+\)hsm inclusion exists in a type~I factor.
\end{proposition}
This is Theorem~\ref{thm:type}. It is a statement about the type of the algebras,
not about observers: a finite-dimensional system has no exact Cut, but an observer
of finite capacity can still couple to a system described by a type III\(_1\)
algebra. Finite instruments and exact Cuts therefore require different
descriptions; Section~\ref{sec:limits} gives their task-level interface.

\begin{proposition}[Accessible algebras and geometry]\label{prop:dictionary}
Let \(A\subset M\) be a subalgebra and
\(N_A^-=\big(\bigcup_{s\le0}\sigma_s(A)\big)''\subset M\). Assume that \(\Omega\) is
cyclic for \(N_A^-\) and \(N_A^-\ne M\). Then:
\begin{enumerate}[label=(\alph*),leftmargin=2em,itemsep=1pt]
\item \((N_A^-\subset M,\Omega)\) is a Cut, with translation group \(U_A\) and
\(N_A^-=U_A(1)MU_A(1)^*\);
\item \(\Delta_{N_A^-}^{it}=U_A(1)\Delta_M^{it}U_A(1)^*\), and \(U_A\) is recovered
from the two modular operators as \(P=\overline{(\log\Delta_{N_A^-}-\log\Delta_M)/2\pi}\);
\item if the inclusion is standard, there is a unique strongly additive local
conformal net \(\mathcal A_A\) with \(\mathcal A_A(0,\infty)=M\),
\(\mathcal A_A(1,\infty)=N_A^-\), translations \(U_A\) and dilations \(\sigma\);
\item \(\bigvee_{t}\sigma_t(A)=M\);
\item if \(A\subset A'\subset M\) and both satisfy the hypotheses, then
\(N^-_A\subset N^-_{A'}\).
\end{enumerate}
\end{proposition}
\begin{proof}
\(\sigma_s(N_A^-)\subset N_A^-\) for \(s\le0\) by construction, and \(\Omega\) is
separating for \(N_A^-\subset M\), so (a) is Theorem~\ref{thm:AZ}(4), and (b) is
(1),(3). (c) is Theorem~\ref{thm:GLW}. (d): \(\bigvee_t\sigma_t(A)\supset
\bigvee_t\sigma_t(N_A^-)=M\) by Lemma~\ref{lem:orbit}, and the left side is
contained in \(M\). (e) is monotonicity of the generated algebra.
\end{proof}

\begin{corollary}[Geometry and modular data determine each other]\label{cor:GLW}
Within the class of standard Cuts, the geometry (the isomorphism class of the
strongly additive net) and the modular data (the isomorphism class of the
triple \((N\subset M,\Omega)\)) determine each other.
\end{corollary}
This is Theorem~\ref{thm:GLW}. The input is the complete standard triple, not a bare
flow: \(\Delta_M\) and \(U\) together with \(M\) and \(\Omega\) fix \(N\) and
\(\Delta_N\), and conversely the two modular operators fix \(U\). The triple replaces
the list ``net, vacuum, translations, dilations''; it does not explain where
\(M\), \(\Omega\) or the Cut come from.

\begin{remark}[Cuts on one orbit are isomorphic]
The Cuts \(N_a\) and \(N_b\) of one net, \(a,b>0\), are exchanged by
\(\sigma_t\) and give isomorphic triples. Distinguishing geometries therefore needs an
invariant of the triple (for instance the spectrum of the conformal Hamiltonian of the
reconstructed net), not merely a different subalgebra. Foundation~I supplies one in this
net: the vacuum character of the hull determines the stabilizer of the seed up to
conjugacy within that model, so seeds with non-conjugate stabilizers give
different vacuum characters and hence non-isomorphic triples~\cite{FoundationI}.
That invariant classifies Foundation~I's ambient-translation Cuts
$U(1)\operatorname{Hull}(S)U(1)^*\subset\operatorname{Hull}(S)$.
It does not classify the past-orbit inclusion $N^-_{O_f}\subset M$
by identifying its reconstructed $U_A$ with the ambient $U$.
Conjecture~\ref{conj:T1}(b) concerns precisely whether that orbit
is an ordinary translated half-line.
\end{remark}

\begin{remark}[Relative position of regions]
For nested intervals \((x_2,x_3)\subset(x_1,x_4)\) the cross-ratio
\(b=\frac{(x_2-x_1)(x_4-x_3)}{(x_3-x_1)(x_4-x_2)}\in(0,1)\) is the unique
modulus under \(PSL(2,\mathbb R)\); in the normal form the outer interval is
\((0,\infty)\) and the inner one is \((b,1)\). Any statement about the relative position of nested regions in the reconstructed net is a statement about \(b\).
\end{remark}

\section{Selection of the accessible algebra by noise}\label{sec:noise}

\subsection{Setting}

\begin{assumption}\label{ass:G}
\(\mathcal A\) is an irreducible, local, M\"obius covariant net on \(S^1\) with
vacuum \((H,\Omega)\). A compact group \(G\) acts by automorphisms \(\alpha_g\),
preserving each \(\mathcal A(I)\), implemented by a strongly continuous unitary
representation \(u(g)\) with \(u(g)\Omega=\Omega\) that commutes with the M\"obius
(and, when present, the diffeomorphism) representation. We write
\(\omega=\langle\Omega,\cdot\,\Omega\rangle\) and \(M=\mathcal A(0,\infty)\).
\end{assumption}

For a closed subgroup \(H\subset G\) put \(Q_H(x)=\int_H\alpha_h(x)\,dh\) with normalised Haar measure.

\begin{lemma}\label{lem:Q}
\(Q_H\) is a normal, faithful, \(\omega\)-preserving conditional expectation of \(M\)
onto the fixed algebra \(M^H\); it is \(M^H\)-bimodular and commutes with the modular
group and with translations.
\end{lemma}
\begin{proof}
For a vector functional \(\rho=\langle\xi,\cdot\,\eta\rangle\) the map
\(h\mapsto\rho\circ\alpha_h\) is norm continuous in the predual, by strong continuity
of \(u\); every normal functional is a norm limit of finite sums of vector
functionals, so the same holds for all normal \(\rho\). The Bochner integral
\(\int_H\rho\circ\alpha_h\,dh\) is therefore a normal functional, and it equals
\(\rho\circ Q_H\); hence \(Q_H\) is normal. Haar invariance gives
\(Q_H^2=Q_H\) and \(\mathrm{Ran}\,Q_H=M^H\); bimodularity is
\(\alpha_h(n_1xn_2)=n_1\alpha_h(x)n_2\) for \(n_i\in M^H\). Since
\(\omega\circ\alpha_h=\omega\), \(\omega\circ Q_H=\omega\), and if \(Q_H(x^*x)=0\)
then \(\omega(x^*x)=0\), so \(x=0\) (faithfulness). Commutation with \(\sigma\) and
translations follows because \(u(g)\) commutes with their unitary groups.
\end{proof}

\subsection{Noise semigroups}

Fix \(\kappa>0\) and put \(D_t^H=e^{-\kappa t}\,\mathrm{id}+(1-e^{-\kappa t})Q_H\).
This is the average channel of a Poisson process of rate \(\kappa\) in which each
jump applies \(\alpha_h\) with \(h\) Haar-distributed in \(H\): the probability of no
jump in time \(t\) is \(e^{-\kappa t}\), and a product of at least one Haar element is
again Haar.

\begin{proposition}\label{prop:haar}
\(D^H_t\) is a norm-continuous semigroup of normal, unital, completely positive,
\(\omega\)-preserving maps on \(M\), and
\begin{enumerate}[label=(\arabic*),leftmargin=2em,itemsep=1pt]
\item \(\mathrm{Fix}(D^H)=M^H\); \(D_t^H-Q_H=e^{-\kappa t}(\mathrm{id}-Q_H)\), so
\(\|D_t^H-Q_H\|\le2e^{-\kappa t}\);
\item \(H_1\subset H_2\Rightarrow M^{H_2}\subset M^{H_1}\); \(\kappa\) changes only the
rate of approach;
\item if \(M^H\ne M\) and \(0\ne y=y^*\in\ker Q_H\), then for \(0<t<\infty\) and \(q=e^{-\kappa t}\) we have \(D_t(y^2)-D_t(y)^2=(1-q)Q_H(y^2)+(q-q^2)y^2\ge0\) and \(\ne0\), so \(D_t\) is a
non-multiplicative channel although it preserves \(\omega\).
\end{enumerate}
\end{proposition}
\begin{proof}
The semigroup law and complete positivity are those of a convex combination of
\(\mathrm{id}\) and \(Q_H\) with \(Q_H^2=Q_H\). (1): if \(D_t(x)=x\) for some
\(t>0\) then \((1-e^{-\kappa t})(x-Q_Hx)=0\). (2) is clear. (3): choose \(x=x^*\notin M^H\)
and \(y=x-Q_Hx\); then \(Q_Hy=0\), \(D_t(y)=qy\), and
\(D_t(y^2)=qy^2+(1-q)Q_H(y^2)\). Since \(Q_H\) is faithful, \(Q_H(y^2)\ne0\).
\end{proof}

\begin{theorem}[Two fixed rotations suffice]\label{thm:four}
Let \(G=SO(3)\) act as in Assumption~\ref{ass:G}. Let \(g_1,\dots,g_m\in SO(3)\) generate a
subgroup whose closure is \(SO(3)\), let \(p_j>0\), \(\sum_jp_j=1\), and put
\(R=\sum_jp_j\alpha_{g_j}\), \(\tilde D_t=\exp[\kappa t(R-\mathrm{id})]\). Then \(\tilde D\) is a
norm-continuous semigroup of normal, unital, completely positive, \(\omega\)-preserving maps with
\(\mathrm{Fix}(\tilde D)=M^{SO(3)}\). The hypothesis holds, in particular,
(i) for \(m=2\) when \(g_1,g_2\) are rotations about two non-parallel axes by angles whose
ratios to \(2\pi\) are irrational (any \(p_1\in(0,1)\)), and (ii) for the symmetric four-operation
choice \(R=\tfrac14(\alpha_{g_1}+\alpha_{g_1^{-1}}+\alpha_{g_2}+\alpha_{g_2^{-1}})\) built from such a pair.
\end{theorem}
\begin{proof}
The generator is bounded, so \(\mathrm{Fix}(\tilde D)=\mathrm{Fix}(R)\). Write
\(\alpha_j\) for the automorphisms and \(u_j\) for their implementers, which
fix \(\Omega\). Then \(\|(\alpha_j(x)-x)\Omega\|^2=2\|x\Omega\|^2-2\mathrm{Re}\langle\alpha_j(x)\Omega,x\Omega\rangle\),
so
\(\sum_jp_j\|(\alpha_j(x)-x)\Omega\|^2=2\|x\Omega\|^2-2\mathrm{Re}\langle R(x)\Omega,x\Omega\rangle\).
If \(R(x)=x\), the left side vanishes, so \(\alpha_j(x)\Omega=x\Omega\) for each \(j\),
and \(\alpha_j(x)=x\) because \(\Omega\) is separating. The set \(\{g:\alpha_g(x)=x\}=\{g:u(g)x\Omega=x\Omega\}\)
is a closed subgroup (strong continuity of \(u\)) that contains every \(g_j\), hence the closure of the
group they generate. For (i): each \(g_j\) generates a dense subgroup of the circle \(T_j\) of rotations
about its axis \(n_j\), so the closure \(H\) of \(\langle g_1,g_2\rangle\) is a closed, hence Lie, subgroup
containing \(T_1\) and \(T_2\). Its Lie algebra contains \(n_1,n_2\in\mathfrak{so}(3)\cong\mathbb R^3\) (bracket = cross
product), hence \(n_1\times n_2\ne0\); the three vectors span \(\mathbb R^3\), so \(H\) contains the
identity component of \(SO(3)\), which is \(SO(3)\). For (ii) the generated group is the same.
Then \(x\) is fixed by all of \(SO(3)\); the converse inclusion is clear.
\end{proof}
Two is not claimed to be the minimal number of operations in every mechanism.

\begin{theorem}[Non-commuting noises select the intersection]\label{thm:multi}
Let \(H_1,\dots,H_m\subset G\) be closed subgroups, \(\kappa_i>0\),
\(L=\sum_i\kappa_i(Q_{H_i}-\mathrm{id})\), \(E_t=e^{tL}\). The \(Q_{H_i}\) need not
commute. Then \(E_t\) is a norm-continuous semigroup of normal, unital, completely
positive, \(\omega\)-preserving maps, with
\[
\mathrm{Fix}(E)=\bigcap_iM^{H_i}=M^{\overline{\langle H_1,\dots,H_m\rangle}},
\qquad
-\langle x\Omega,Lx\Omega\rangle=\sum_i\kappa_i\|(x-Q_{H_i}x)\Omega\|^2,
\]
and \(E_t\) converges strongly on \(L^2(M,\omega)\) to the Haar average over the
closed group generated.
\end{theorem}
\begin{proof}
With \(\lambda=\sum\kappa_i\) and \(\bar R=\sum\kappa_iQ_{H_i}/\lambda\),
\(E_t=e^{-\lambda t}\sum_n(\lambda t)^n\bar R^n/n!\) is a Poisson mixture of powers of a
state-preserving unital completely positive map. On the GNS space each \(Q_{H_i}\)
is the orthogonal projection \(P_i=\int_{H_i}u(h)\,dh\); hence \(L\) is bounded,
self-adjoint and nonpositive there, and the displayed identity holds. \(Lx=0\) forces
\(Q_{H_i}x=x\) for every \(i\), because \(\Omega\) is separating. The spectral
theorem gives strong convergence to the projection onto \(\ker L=\bigcap_i\mathrm{Ran}P_i\),
which is the Haar average over the closed group generated.
\end{proof}
No uniform spectral gap or operator-norm rate is claimed in infinite dimension.

\subsection{SU(2) level one and the U(1) contrast}

Let \(\mathcal A_{SU(2)_1}\) be the vacuum net of the loop group \(LSU(2)\) at level
one; \(SO(3)\) acts properly by the adjoint action, and \(\mathrm{Vir}_1\) denotes the
Virasoro subnet of the Sugawara stress tensor, \(c=1\).

\begin{theorem}[Rehren {\cite{Rehren}}]\label{thm:Rehren}
\(\mathcal A_{SU(2)_1}(I)^{SO(3)}=\mathrm{Vir}_1(I)\) for every interval \(I\).
\end{theorem}
The proof in \cite{Rehren} decomposes the vacuum representation by isospin: its
character is \(\sum_j\chi_j(z)\,t^{j^2}(1-t^{2j+1})\prod_n(1-t^n)^{-1}\), the
\(SO(3)\)-invariant part is the \(j=0\) term, which is the character of the
irreducible Virasoro vacuum representation at \(c=1\) (Segal \cite{Segal}, Kac--Raina \cite{KR}), so the
invariant vacuum space equals the Virasoro vacuum space. Haag duality of the Virasoro net \cite{BSM} and locality then give the chain
\(\mathrm{Vir}(I)\subset\mathcal A(I)^G\subset(\mathcal A(I^c)^G)'\subset\mathrm{Vir}(I^c)'=\mathrm{Vir}(I)\)
on that space, and local faithfulness of the representation lifts the equality to
\(\mathcal A\). We recomputed the Virasoro vacuum module in the charge-zero Fock space
with exact rational arithmetic (Appendix~\ref{app:fock}): its dimensions at levels \(0,\dots,9\)
are \(1,0,1,1,2,2,4,4,7,8=p(n)-p(n-1)\), as the character requires. Since
\(\mathrm{Vir}_1\) is strongly additive, the reconstruction of Section~\ref{sec:cut}
returns \(\mathrm{Vir}_1\) itself and not a dual net.

\begin{corollary}\label{cor:selected}
In \(\mathcal A_{SU(2)_1}\) the exact fixed algebra of the Haar-jump noise
\(D^{SO(3)}\), of the semigroup \(\tilde D\) built from finitely many fixed rotations that generate \(SO(3)\) (two suffice), and of the pair of
non-commuting axis noises \(H_z,H_x\) is \(\mathrm{Vir}_1(0,\infty)\), a proper
type III\(_1\) subfactor of \(M\).
\end{corollary}
Properness: \(\alpha\) is non-trivial on \(\mathcal A\) and trivial on \(\mathrm{Vir}_1\).

\begin{proposition}[The \(\mathbb Z_2\) noise of the \(U(1)\) current does not select the Virasoro net]\label{prop:U1}
Let \(\mathcal A_{U(1)}\) be the net of the \(U(1)\) current \(J\), \(\gamma:J\mapsto-J\),
\(F=\tfrac12(\mathrm{id}+\gamma)\). Then
\(\mathrm{Vir}_1(I)\subsetneq\mathcal A_{U(1)}(I)^{\mathbb Z_2}\subsetneq\mathcal A_{U(1)}(I)\).
\end{proposition}
\begin{proof}
At level 4 the even Fock space is spanned by \(J_{-1}^4\Omega,J_{-3}J_{-1}\Omega,J_{-2}^2\Omega\),
dimension 3, whereas the Virasoro vacuum module at level 4 has dimension 2, spanned by
\(L_{-4}\Omega=(J_{-3}J_{-1}+\tfrac12J_{-2}^2)\Omega\) and
\(L_{-2}^2\Omega=(\tfrac14J_{-1}^4+J_{-3}J_{-1})\Omega\), with Gram matrix
\(\left(\begin{smallmatrix}5&3\\3&9/2\end{smallmatrix}\right)\). The vector
\(u=\tfrac12J_{-1}^4\Omega-J_{-3}J_{-1}\Omega+\tfrac34J_{-2}^2\Omega\) is even,
orthogonal to both, satisfies \(L_mu=0\) for \(m=1,\dots,4\), and has
\(\|u\|^2=27/2\) (Appendix~\ref{app:fock}). If \(\mathrm{Vir}_1(I)=\mathcal A(I)^{\mathbb Z_2}\), then
by Reeh--Schlieder their vacuum-cyclic subspaces coincide: \(F(\mathcal A(I))\Omega\)
is dense in the even subspace and \(\mathrm{Vir}_1(I)\Omega\) is dense in the
Virasoro vacuum space; but \(u\) lies in the former and not the latter. The second
inclusion is strict because \(\gamma\) is non-trivial.
\end{proof}
So the \(\mathbb Z_2\) noise of the \(U(1)\) current protects the stress tensor (and the
whole even algebra) but does not select \(\mathrm{Vir}_1\); a noise that selects
\(\mathrm{Vir}_1\) in that net must be specified differently, for instance by the
conditional expectation onto \(\mathrm{Vir}_1\) itself. The \(SU(2)\) construction
has the advantage that the noise is specified by internal rotations, not by the target
algebra.

\begin{remark}[Fixed algebra, attracting state, correctable algebra]
The three notions are of different types. The semigroups above preserve the vacuum
state and act as the identity on the fixed algebra, so information in the fixed algebra
is passively preserved and can be recovered by the identity map. A reset channel \(E_t(A)=e^{-\gamma t}A+(1-e^{-\gamma t})\mathrm{tr}(\rho_*A)\mathbf 1\)
has an attracting state \(\rho_*\) but trivial fixed algebra, and actively protected
information is described by correctable algebras (B\'eny--Kempf--Kribs \cite{BKK}). We do
not identify a fixed algebra with an attracting state, and we do not claim that the full
correctable structure of these semigroups equals their fixed algebra; exact fixed
algebras are, moreover, discontinuous in the dynamics (Section~\ref{sec:limits}).
\end{remark}

\section{Access: one observable generates the selected algebra}\label{sec:access}

Let \(\mathcal A\) be diffeomorphism covariant with \(c>0\), \(M=\mathcal A(0,\infty)\),
\(B=\mathrm{Vir}_c(0,\infty)\). For \(f\in C_c^\infty((0,\infty),\mathbb R)\), \(f\ne0\), the
\emph{seed} is \(O_f=\tfrac2\pi\arctan T(f)\in B\), a bounded self-adjoint
observable. An \emph{admissible} algebra is a unital von Neumann algebra
\(N\subset M\) with \(\sigma_s(N)=N\) for all \(s\in\mathbb R\) and \(T_a(N)\subset N\)
for \(a\ge0\). The \emph{hull} \(\mathrm{Hull}_{\sigma,T}(O_f)\) is the intersection of all
admissible \(N\) containing \(O_f\).

\begin{theorem}[Generation]\label{thm:hull}
\(\mathrm{Hull}_{\sigma,T}(O_f)=\mathrm{Vir}_c(0,\infty)\) for every nonzero
\(f\in C_c^\infty((0,\infty),\mathbb R)\).
\end{theorem}
Appendix~\ref{app:hull} proves this for arbitrary nonzero \(f\). The stronger
modular-only result below applies to a specified bump class, not to all \(f\).

\begin{proposition}[Forward translations alone do not suffice]\label{prop:transonly}
If \(\mathrm{supp}f\subset(a,b)\) with \(a>0\), then the smallest algebra containing
\(O_f\) and stable under \(T_s\), \(s\ge0\), lies in \(\mathrm{Vir}_c(a,\infty)\subsetneq B\).
\end{proposition}
\begin{proof}
All forward translates of \(O_f\) lie in \(\mathrm{Vir}_c(a,\infty)\), which is stable
under forward translations. If \(\mathrm{Vir}_c(a,\infty)=B\) then
\(\mathrm{Vir}_c(0,a)\subset B\cap B'\) by locality, which is trivial as \(B\) is a
factor; but \(\mathrm{Vir}_c(0,a)\) is not trivial.
\end{proof}
For the logarithmic bump class of Theorem~\ref{thm:singleorbit}, dilations
alone suffice by Corollary~\ref{cor:modonly}. The unrestricted stress-seed question
remains open. Foundation~I gives a separate modular-only generation theorem for
nonzero smeared current seeds~\cite{FoundationI}.

A seed here denotes an ideal observable, not a grant of every function of it
as an actual measurement or gate. An injective bounded function has the
same spectral algebra as its generator; its implementation has a separate price.
Foundation~I distinguishes an effect, its spectral hull and the full physical
instrument. HAFF-E supplies a concrete contrast: a partially polarized memory
produces an unsharp effect generating the same pointer algebra at every nonzero
polarization, while its one-read error varies continuously~\cite{PaperE}.
The same distinction applies to \(O_f\) and to the ideal flows used below.

\begin{example}[Exact fixedness does not freeze the modular flow]
In \(M_2\bar\otimes M_2\) with state \(\mathrm{diag}(.8,.2)\otimes\tfrac12 I\), let the
noise be \(\mathrm{id}\otimes\) depolarisation and the observable \(X\otimes1\). The
observable is fixed by the noise, but \(\sigma_t(X\otimes1)=(\cos(t\ln4)X-\sin(t\ln4)Y)\otimes1\)
generates the non-commutative algebra \(M_2\otimes1\). Exact invariance under the noise
does not by itself prevent the modular flow from moving a fixed observable through a
non-commutative algebra.
\end{example}

\begin{theorem}[Access closure equals noise-selected algebra]\label{thm:main}
In \(\mathcal A_{SU(2)_1}\), let \(O_f\) be built from the Sugawara stress tensor. Then
\[
\mathrm{Hull}_{\sigma,T}(O_f)=\mathrm{Vir}_1(0,\infty)=M^{SO(3)}=\mathrm{Fix}(D^{SO(3)})=\mathrm{Fix}(\tilde D).
\]
\end{theorem}
\begin{proof}
The first equality is Theorem~\ref{thm:hull} applied inside \(M\): the argument
uses only unitaries \(e^{iuT(h)}\in B\subset M\). The second is Theorem~\ref{thm:Rehren}
(the half-line is a M\"obius image of an interval), and the last two are
Corollary~\ref{cor:selected}.
\end{proof}
Two different structural definitions meet here: the algebra generated from one
observable by dilations and forward translations, and the algebra left invariant by an
explicit random-rotation noise. The equality is specific to this net and group; it
fails for the \(\mathbb Z_2\) noise of the \(U(1)\) current
(Proposition~\ref{prop:U1}).

\begin{remark}[Two different inclusions]\label{rem:twolevel}
The algebra selected by noise is \(B=M^{SO(3)}=\mathrm{Vir}_1(0,\infty)\subset
M=\mathcal A_{SU(2)_1}(0,\infty)\). By Lemma~\ref{lem:Q} there is an
\(\omega\)-preserving faithful normal conditional expectation of \(M\) onto \(B\), so by
Takesaki's theorem \cite{Takesaki} \(\sigma^M_s(B)=B\) for all \(s\in\mathbb R\) and
\(\sigma^M|_B\) is the modular group of \(B\) in its own vacuum representation, so stability of \(B\)
under \(\sigma\) in one direction holds trivially. What fails is cyclicity. For \(x\in M\),
\(\int u(g)x\Omega\,dg=Q_{SO(3)}(x)\Omega\) because \(u(g)\Omega=\Omega\); so the projection onto the
\(SO(3)\)-invariant subspace \(H^{SO(3)}\) maps \(M\Omega\) onto \(B\Omega\), and since \(M\Omega\) is dense,
\(\overline{B\Omega}=H^{SO(3)}\), a proper subspace because the vacuum representation of
\(\mathcal A_{SU(2)_1}\) contains non-trivial isospin sectors. \(\Omega\) is therefore not cyclic for
\(B\) on the ambient Hilbert space, and \(B\subset M\) is not a Cut in the sense of
Definition~\ref{def:hsm} (which asks \(\Omega\) to be cyclic and separating for the subalgebra; a proper
subalgebra that is \(\sigma\)-invariant and has \(\Omega\) cyclic would contradict
Lemma~\ref{lem:nest}). The Cuts of the geometric construction are the half-line subalgebras
\(N_a=\mathrm{Vir}_1(a,\infty)\subset B\), taken in the vacuum-cyclic representation
of \(B\) (a factor, so the restriction to that subspace is faithful). The two structures
meet in one technical model: noise fixes \(B\), and \(B\) carries the net whose Cuts
are the \(N_a\). The noise selects \(B\) without producing its vacuum or modular dynamics.
Theorem~\ref{thm:singleorbit}, applied in \(H^{SO(3)}\), then constructs a
standard Cut from a local seed and that modular dynamics, and reconstructs its
translations. These need not be assumed as a separate generation permission.
The underlying Virasoro carrier and state are still inherited from the
ambient model.
\end{remark}

\subsection{A Cut generated by one past modular orbit}\label{sec:singleorbit}

Here \(M=\mathrm{Vir}_c(0,\infty)\) acts on its own vacuum-cyclic space.
Write
\[
 N^-_{O_f}=W^*(\sigma_s(O_f):s\le0),\qquad
 H_\sigma(O_f)=W^*(\sigma_s(O_f):s\in\mathbb R).
\]
The first algebra is automatically half-sided invariant; it still needs
cyclicity and properness before it is a Cut.

\begin{lemma}[Smooth flow closure]\label{lem:flowclosure}
For a unital von Neumann algebra \(N\), the real compactly supported vector
fields whose stress flows \(e^{itT(v)}\) belong to \(N\) for every \(t\in\mathbb R\)
are closed under real linear combinations, under \(C^\infty\) limits with
supports in a fixed compact set, and under the ordinary vector-field bracket
\([v,w]=vw'-v'w\).
\end{lemma}
The decisive bracket step is geometric conjugation of stress flows,
followed by a fixed-compact smooth limit; scalar projective phases do
not affect membership in a unital algebra. The full argument is in
Appendix~\ref{app:singleorbit-proofs}.


\begin{theorem}[A single-orbit standard Cut]\label{thm:singleorbit}
Choose \(0<a<b\), \(\alpha=\log a\), \(\beta=\log b\), and a real
\(g\in C_c^\infty(\mathbb R)\) with support \([\alpha,\beta]\).
Suppose \(g>0\) on \((\alpha,\beta)\), and \(g\) is strictly increasing and
strictly decreasing on some interior neighborhoods of its left and right
endpoints, respectively. Put \(f(x)=xg(\log x)\) and
\(O_f=(2/\pi)\arctan T(f)\). Then
\begin{equation}\label{eq:singleorbit}
 \mathrm{Vir}_c(b,\infty)\subset N^-_{O_f}
 \subset\mathrm{Vir}_c(a,\infty)\subsetneq M .
\end{equation}
The triple \((N^-_{O_f}\subset M,\Omega)\) is a proper standard \(+\)hsm
inclusion. Its translations and strongly additive net are recovered from its
modular operators.
\end{theorem}
The proof constructs a positive bracket kernel near the right
endpoint of the logarithmic bump. Its permitted right translates
approximate every compact stress field on the tail \((b,\infty)\).
Locality and Reeh--Schlieder then give properness, cyclicity and
standardness. Appendix~\ref{app:singleorbit-proofs} gives the kernel
construction and all closure steps.


\begin{corollary}[Modular flow alone generates the full half-line]\label{cor:modonly}
For the seeds in Theorem~\ref{thm:singleorbit},
\[
 H_\sigma(O_f)=\mathrm{Vir}_c(0,\infty).
\]
\end{corollary}
\begin{proof}
The two-sided hull contains \(N^-_{O_f}\), hence
\(\mathrm{Vir}_c(b,\infty)\), and is invariant under every dilation. It contains
\(\mathrm{Vir}_c(\lambda b,\infty)\) for every \(\lambda>0\).
Every compact stress field on \((0,\infty)\) lies in one of these half-lines,
so their generated algebra is \(M\).
\end{proof}

For example, \(g(u)=\exp[-1/((u-\alpha)(\beta-u))]\) inside its support and
zero outside satisfies the hypotheses. Positivity can be reversed globally
without changing the spectral algebra. The theorem is not an all-seed
classification: sign-changing and more irregular endpoint behavior are not
covered by the positive-kernel argument.

The lower inclusion is enough to prove a Cut, but does not identify
\(N^-_{O_f}\) with, or distinguish it from, an ordinary translated half-line.
For \(c>1\), the standard reconstruction is strongly additive even though the
original Virasoro net is not; one must not infer that the reconstructed full
net equals the original net from the shared half-line algebra alone.
For \(c=1\) the carrier is strongly additive, but whether this seed changes the
isomorphism class of the reconstructed triple still requires an invariant.

\begin{proposition}[Cyclicity supplies the remaining Cut conditions]\label{prop:standard}
For arbitrary nonzero \(f\) supported in a compact subset of \((a,b)\),
\(a>0\), in the same vacuum representation, cyclicity of \(N^-_{O_f}\)
implies that \(N^-_{O_f}\subset M\) is a proper standard \(+\)hsm inclusion.
\end{proposition}
\begin{proof}
The past orbit lies in \(\mathrm{Vir}_c(a,\infty)\subsetneq M\), and
\(\mathrm{Vir}_c(0,a)\subset M\cap(N^-_{O_f})'\). Locality and
Reeh--Schlieder give standardness, while separatingness is inherited from \(M\).
\end{proof}

\begin{conjecture}[General seeds and the position of the Cut]\label{conj:T1}
(a) For every nonzero real compact stress seed on \((0,\infty)\), the past
modular orbit is vacuum-cyclic.
(b) Whenever the past orbit is cyclic, \(N^-_{O_f}\) is not
\(\mathrm{Vir}_c(a',\infty)\) for any \(a'>0\).
\end{conjecture}
Part (a) is proved above for the endpoint-monotone positive log-bump class.
Beyond that class, bracket closure remains available but a localized kernel of
controlled nonzero mass has not been obtained in general.
Part (b) is independent of cyclicity. Near the left support endpoint, finite
bracket expressions built from right translates are multiples of the initial
field. This germ restriction does not survive automatically through strong
operator closure, so it cannot prove (b). A non-isomorphic geometry would
require a further invariant of the full standard triple.
The positive construction narrows the formation problem without assuming its
remaining general conclusions.

\section{Response of the selected algebra to excitation}\label{sec:response}

Let \(f\in C_c^\infty((a_0,\infty),\mathbb R)\), \(a_0>0\), \(\rho=\mathrm{Exp}f\),
\(w=\rho''/\rho'\), and \(0\le a<a_0\). Put \(N_a=\mathrm{Vir}_c(a,\infty)\).

\begin{theorem}[Cut displacement]\label{thm:displacement}
With \(U=V(\rho)\), \(\omega_\pm\) as in Section~\ref{sec:hollands},
\[
D_a:=S_{N_a}(\omega\|\omega_+)=S_{N_a}(\omega_-\|\omega)=\frac c{24}\int(x-a)\,w(x)^2\,dx .
\]
Consequently
\[
-\partial_aD_a=2\pi E,\qquad D_0-D_a=2\pi aE,\qquad
E=\frac c{48\pi}\int w^2\,dx=\int\langle T(x)\rangle_{\omega_-}dx,
\]
where \(\langle T\rangle_{\omega_-}=-\tfrac c{24\pi}\mathrm{Sch}(\rho)\) is the stress
expectation in \(\omega_-\).
\end{theorem}
\begin{proof}
Since \(\omega_+\circ\mathrm{Ad}\,U=\omega\) and \(\omega\circ\mathrm{Ad}\,U=\omega_-\),
unitary covariance of the relative entropy gives
\(S(\omega\|\omega_+)=S(\omega_-\|\omega)\) on any algebra containing \(U\), in
particular \(N_a\). Translation covariance moves the algebra \(N_a\) to \(N_0\) and
\(\rho\) to \(\tilde\rho(y)=\rho(y+a)-a\), whose generating field
\(\tilde f(y)=f(y+a)\) is supported in \(y>a_0-a>0\); \eqref{eq:hollands} applied to
\(\tilde\rho\) gives \(\frac c{24}\int y\,(\partial\log\tilde\rho')^2dy=\frac c{24}\int(x-a)w^2dx\).
The same value is given by the exact identity for finite-amplitude local unitaries (not the linearised first law),
\(S_{N_a}(\omega_-\|\omega)=\langle U^*\Omega,-\log\Delta_{N_a}U^*\Omega\rangle
=2\pi\int(x-a)\langle T\rangle_{\omega_-}dx\), obtained from the transformation law of
\(T\) and \(\int(x-a)\mathrm{Sch}(\rho)\,dx=-\tfrac12\int(x-a)w^2dx\). Differentiating in
\(a\) gives the remaining formulas.
\end{proof}

\begin{corollary}\label{cor:hessian}
\(G_a(f)=\partial^2_\theta D_a(\mathrm{Exp}\,\theta f)|_{\theta=0}=\frac c{12}\int(x-a)f''(x)^2dx\), which is \(>0\) for every non-zero test field \(f\),
and \(G_0-G_a=\tfrac{ca}{12}\int f''^2\).
\end{corollary}

\begin{proposition}[A displaced boundary is not reversible recoding]
\label{prop:access-recoding}
In Theorem~\ref{thm:displacement}, let \(0<a<b<a_0\) and \(E>0\).
Write \(\omega_a=\omega|_{N_a}\) and
\(\varphi_a=\omega_-|_{N_a}\). Restriction from \(N_a\) to \(N_b\)
loses
\[
 S_{N_a}(\varphi_a\|\omega_a)-S_{N_b}(\varphi_b\|\omega_b)
       =2\pi(b-a)E>0 .
\]
There is no normal unital completely positive map
\(\mathcal R:N_a\to N_b\) with
\(\omega_b\circ\mathcal R=\omega_a\) and
\(\varphi_b\circ\mathcal R=\varphi_a\).
Thus no common physical recovery restores this pair after restriction,
although the vacuum Cut triples at \(a\) and \(b\) are isomorphic.
\end{proposition}
\begin{proof}
The entropy difference is Theorem~\ref{thm:displacement}.
If such a recovery existed, monotonicity under \(\mathcal R\)
would give
\(S_{N_a}(\varphi_a\|\omega_a)\le
  S_{N_b}(\varphi_b\|\omega_b)\), contradicting strict loss.
The two vacuum triples are related by a dilation, as noted after
Corollary~\ref{cor:GLW}; that isomorphism also transports states.
Here the excitation is held fixed rather than transported with the Cut.
\end{proof}

A simultaneous relabelling of the algebra, states and instruments
preserves their complete experiment. A boundary moved relative to a
fixed excitation need not do so. The proposition establishes this
difference for an ideal state pair. A finite test must still specify
the permitted instruments; it is not a finite-sample estimate of Araki
relative entropy merely because that entropy is finite.

\begin{proposition}[Noise does not reduce the response]\label{prop:Fprime}
Let \(E:M\to B\) be a normal conditional expectation and \(\iota:B\to M\) the
inclusion. If two normal states \(\varphi,\psi\) on \(M\) factor through \(E\)
(\(\varphi=\varphi|_B\circ E\), \(\psi=\psi|_B\circ E\)), then
\(S_M(\varphi\|\psi)=S_B(\varphi|_B\|\psi|_B)\).
\end{proposition}
\begin{proof}
Monotonicity under the restriction \(\iota\) gives \(S_B\le S_M\); monotonicity under
the unital completely positive map \(E\) gives \(S_M(\varphi\|\psi)\le S_B\).
\end{proof}
In \(\mathcal A_{SU(2)_1}\), for \(U\in\mathrm{Vir}_1\), the states \(\omega\) and
\(\omega_\pm\) factor through \(Q_{SO(3)}\): put \(W_+=U\), \(W_-=U^*\), so that
\(\omega_\pm(x)=\omega(W_\pm^*xW_\pm)\) with \(W_\pm\in M^{SO(3)}\); since \(Q_{SO(3)}\) is
\(M^{SO(3)}\)-bimodular and \(\omega\circ Q_{SO(3)}=\omega\),
\(\omega_\pm(Q_{SO(3)}x)=\omega(Q_{SO(3)}(W_\pm^*xW_\pm))=\omega_\pm(x)\); hence the information
carried by an excitation of the selected algebra is the same in the full accessible
algebra \(M\) as in \(\mathrm{Vir}_1\); the noise does not degrade it. This statement is
about states generated inside the selected algebra, not about arbitrary information.

\begin{remark}[What the identity is not]
The identity is a statement about relative entropy on a family of algebras, and
\(-\partial_aD_a\) is a null-energy integral. It is not a Clausius relation for a
thermodynamic entropy, and it does not produce an area or a curvature: moving a
region in flat vacuum changes the algebra and its modular generator without creating
curvature. The Hessian of the relative entropy corresponds to the Bogoliubov--Kubo--Mori
metric, not to the symmetric-logarithmic-derivative quantum Fisher information. The
structure is that of the half-sided modular inclusions used by Ceyhan and Faulkner
\cite{CF} to relate shape derivatives of relative entropy to null energy, and null-plane
algebras carry a copy of such a structure for each transverse position \cite{CTT}; a
transverse family with coupling is where area and focusing could enter, and is not
treated here.
\end{remark}

\section{Limits of the results}\label{sec:limits}

\paragraph{Exact algebras and finite tasks.} Exact fixed algebras and hulls
are discontinuous in the dynamics. Two qubits with local \(Z\)-dephasing at
rate \(\kappa\) and coupling \(\epsilon X_1X_2\) have the fixed algebra
\(\mathrm{span}\{1,Z_1Z_2\}\) for every \(\epsilon\ne0\), versus a
four-dimensional diagonal algebra at \(\epsilon=0\).
For \(|\epsilon|<\kappa\), the slowest nonzero rate is
\(2(\kappa-\sqrt{\kappa^2-\epsilon^2})\), so four slow modes persist on a
window of order \(\kappa/\epsilon^2\). Their span is a rotated commutative
algebra, but finite discrimination errors, not its exact fixedness, determine
a reader's available tasks. HAFF-B treats such windows and the continuity of
finite task risk~\cite{PaperB}.

Nor does the new generation theorem give a cheap continuum limit.
For the normalized bump \(g(u)=e^{4-1/(u(1-u))}\), the bracket mass at
\(\varepsilon=.05\) is about \(1.1\times10^{-31}\). Dividing by that mass and
taking spectral, Trotter and smooth limits are legitimate algebraic steps;
they supply no bounded-control or finite-time certificate.

There is nevertheless a useful finite consequence of the constructive proof.
Its scope is conditional on implementing the elementary flows, not on merely
measuring the bounded seed.

\begin{proposition}[A finite task consequence of flow generation]\label{prop:finitewords}
Take \(f,g,a,b\) as in Theorem~\ref{thm:singleorbit}, with
\(f_c(x)=xg(\log x-c)\).
Suppose the elementary gates \(e^{itT(f_c)}\), for finitely specified real
\(t\) and \(c\ge0\), can be implemented to arbitrary specified error.
For every compact stress field \(h\) supported in \((b,\infty)\), every
joint density operator \(\rho\) on the vacuum Hilbert space and a retained reference, and
every \(\delta>0\), a finite word of such gates and their inverses approximates
the channel of \(e^{iT(h)}\) on \(\rho\) within trace distance \(\delta\).
The statement holds uniformly over a specified finite list of joint states.
If slot \(j\) is implemented within half-diamond error \(\epsilon_j\), the
additional trace-distance error is at most \(\sum_j\epsilon_j\).
\end{proposition}
\begin{proof}
The proof of Theorem~\ref{thm:singleorbit} uses conjugations, real linear
combinations implemented by Trotter products, and fixed-compact smooth limits.
For each constructed flow these give iterated strong limits of phase-adjusted
finite words in the elementary flow unitaries. Scalar phases do not change
the resulting channels. Their adjoints converge strongly too, since the
limits are unitary and the inverse-time construction applies.
On any finite set of test vectors, successive approximations can be chosen
in the order in which these limits occur to produce one finite word with
arbitrarily small vector error. On the separable vacuum Hilbert space a
diagonal choice gives strong convergence of a sequence of words to
\(e^{iT(h)}\).

Strong convergence of uniformly bounded unitaries implies trace-norm
convergence of conjugation on every trace-class joint state, first on
finite-rank operators on the field and reference and then by trace-norm
approximation. The same choice works for a finite list of states.
For imperfect gates, telescoping through the same circuit and contractivity
of channels gives the sum of slot errors, including the retained reference.
\end{proof}

This proposition is not uniform over all input states and gives no rate or
optimal cost. Its targets are the stress flows constructed above, not every
unitary or channel in the generated von Neumann algebra. Each finite
approximation has a finite word length, gate
precision, maximum modular shift and preparation budget that must be charged.
\(W^*(O_f)\) alone does not license its spectral functions as hardware.
In a neutral current realization, that permission has a concrete
implementation rather than an independent quadratic-control assumption.

\begin{corollary}[Finite flows from local current control]
\label{cor:physical-finite-stress}
Use the specified circle realization with neutral current Hamiltonian
\(H=L_0\), restricted at input to its \(c=1\) Virasoro vacuum subspace.
Suppose signed local linear-current coupling, degenerate-helper
rotations, record phase gates and calibrated clock preparation are
available during coherent control intervals. Then the elementary-flow
permission of Proposition~\ref{prop:finitewords} can be implemented
by a finite positive autonomous program. For every fixed target there
and \(E<\infty\), its joint error can be made arbitrarily small,
uniformly over inputs with \(\operatorname{tr}(H\rho)\le E\), retaining
the old references, helper and original clock. The conclusion also
holds on a specified finite list of joint states without a uniform
energy bound.
\end{corollary}
\begin{proof}
Q-RAIF~C's current-stress compiler~\cite{QRC} selects finite current
quadratures using full-period linear pulses, synthesizes their
squares with a retained Pauli helper, and proves convergence of the
current-factor cutoff to the Sugawara stress tensor. Its positive
single-clock realization includes free propagation and both signs
of each gate. Map the finitely many compact line smearings to the
specified circle chart, accounting for scalar coordinate phases.
On \(P_M=1_{[0,M]}(H)\), strong finite-word convergence becomes
operator-norm convergence because \(P_M\) has finite rank.
Approximate the finite set of prefix vectors when compiling the
word and telescope the errors. For a joint mean-energy input the
discarded proof tail costs at most \(2\sqrt{E/(M+1)}\), independently
of reference dimension. For a finite list of arbitrary joint states,
their spectral tails instead tend to zero by trace-class convergence.
The proof uses these projections only to bound errors, not as physical
measurements or postselection.
\end{proof}

The physical choice has a bounded future test. Use the Fourier convention
\[
 f_n=(2\pi)^{-1}\int_0^{2\pi} f(\theta)e^{-in\theta}d\theta.
\]
Take real compact \(f\) with \(b=\operatorname{Im}f_2\ne0\), and
\(q_1=(a_1+a_1^*)/\sqrt2\). The vacuum flow satisfies
\[
 \langle e^{itT(f)}\Omega,q_1^2e^{itT(f)}\Omega\rangle
       =\tfrac12-bt+O(t^2).
\]
Only the first-mode pair term contributes to this derivative.
A later current Ramsey test of \(e^{i\lambda q_1}\) therefore has
\(p_+'(0)=b\lambda^2e^{-\lambda^2/4}/4\ne0\).
A retained binary record selecting the two signs changes this later
readout for sufficiently small nonzero \(t\) and implementation error.
The test requires another paid linear pulse, helper and measurement.
Intermediate current pulses can leave the Virasoro vacuum subspace;
their endpoint leakage is included in the joint error. They are not
protected by the always-on \(SO(3)\) noise merely because the target
stress flow is invariant. The carrier, coherent-control interval and
growing precision budgets remain inputs; the finite program does not
establish the exact Cut limit or its maintenance.

Foundation~I gives a sharper, explicit bridge for finite pointer tasks:
fresh partially polarized records approximate a sharp pointer instrument in
half-diamond distance by their exact majority error, and errors add through
finite adaptive circuits~\cite{FoundationI}. HAFF-A supplies a different
finite achievement: actual calibrated fragment readouts certify distinct
correlation-graph geometries inside a stated GHZ protocol~\cite{PaperA}.
Neither result asserts that a finite apparatus has become an exact type-III Cut.

Sharing a carrier algebra does not fix its state: the vacuum and a
local Virasoro excitation are both fixed by $D^{SO(3)}$ but have
$D_0>0$. Conversely the two tensor-factor algebras intersect only
in scalars while the correlated state
$(1-\epsilon)(|00\rangle\langle00|+|11\rangle\langle11|)/2+\epsilon I/4$
has $Z$ disagreement $\epsilon/2$. Foundation~II treats this distinction.

\section{Geometry, access and the organization of a perspective}\label{sec:hard}

The positive claim is that a perspective is constrained by relations
within the world it encounters. Its boundary determines which
distinctions can be used; a retained model of those conditions can
change what the process next reads and does. Geometry, access and
self-description then meet through physical operations. They are not
independent labels attached after those operations have been specified.

\subsection{What geometry fixes, and what it leaves open}

For a standard Cut, the complete modular triple and the reconstructed
net determine each other. This rules out their arbitrary independent
variation within that class. It does not choose an ontological priority
between two equivalent specifications. Nor does it identify the net's
modular data with all the information in a viewpoint.

Proposition~\ref{prop:access-recoding} sharpens the distinction.
Dilations relate the vacuum Cut triples at \(a\) and \(b\), so their
reconstructed geometric isomorphism class is the same. Holding an
excitation fixed while restricting access from \(N_a\) to \(N_b\)
nevertheless loses information that no common channel can restore for
the vacuum and excited state. The source--boundary relation is physical
data beyond the abstract isomorphism class. Recoding alone cannot
restore the lost state-pair information.

There are consequently two ways in which structure enters a perspective.
The chosen boundary and state constrain its geometric description.
The boundary's embedding relative to actual sources and instruments
constrains its usable distinctions. The first is exact in the standard
Cut correspondence; the second is witnessed by the entropy response
and by finite instrument laws. The single-orbit theorem removes a
separate Cut input for its seed class, but it uses the carrier and
vacuum dynamics. These are formation results with stated inputs,
rather than a universe waiting for a person to observe it.

Finite central records do not supply the same object. In Foundation~I's
stationary commuting-noise model the modular flow fixes the central
projections, so their stable hull is their abelian algebra and produces
no proper modular Cut \cite{FoundationI}. The chiral construction here
starts from a net and vacuum and generates a one-sided orbit inside
the appropriate vacuum representation. Neither construction is a
missing limit theorem connecting the two.

\subsection{What history must change}

Functional self-reading means that an internally available, revisable
prediction of the system's own observation or action map is updated
from records and alters subsequent reading or action
\cite{WorldImprint}. An accurate external account of the same map does
not yet satisfy that causal role. The distinction concerns the return
from acquired information to this process, not a special format in
which it has to describe itself.

Consciousness~III gives an exact way to test the role. Let \(R\) be
the acquired record, \(Z=f(R)\) what is retained, and \(\mathcal U\)
the allowed finite continuation menu, including acquisition and
actuation prices. For bounded task losses, let \(L_u(R)\) be the
conditional loss of action \(u\). Then
the decision price of keeping only \(Z\) is
\begin{equation}
 \Delta_{\rm dec}
 =\mathbb E\min_{u\in\mathcal U}\mathbb E[L_u(R)\mid Z]
   -\mathbb E\min_{u\in\mathcal U}L_u(R)\ge0 .
 \label{eq:perspective-decision}
\end{equation}
It vanishes precisely when one action is optimal almost surely in
each retained-record cell \cite{PaperCIII}. That paper proves the
identity and its equality condition; it is used here as an interface,
not a second model or a new uncertainty law.

A confidence score need not be such a sufficient state. In its
world--sensor--record-frame task, the two acquired posteriors
\((1,9,1)/11\) and \((1,1,9)/11\) both have confidence \(9/11\).
Retaining which internal port was queried permits restoration with
failure \(2/11\); erasing the port and every correlated record while
retaining confidence alone gives failure \(6/11\). The loss difference
is \(4/11\), conditional on the declared event and continuation menu.
This is a failure of a particular compressed self-description, not an
additional thermodynamic or relative-entropy term.

The result explains why a maintained perspective needs more than a
confidence display. It must preserve the distinctions that organize its
possible corrections. Conversely, redundant logs can be lost without
harm if the sufficient causal state and update program survive.
A faithful copy of that state into replacement hardware can preserve
the permitted future laws. A common relabelling also preserves them.
History has no unobservable privilege after all of its causally
relevant organization has been reproduced.

The relation to geometry is through actual access. History can change
an instrument, the state on which it acts, or which operations become
available; it can also merely recode an unchanged experiment. The
first kind of change can alter the world available to the process,
as Proposition~\ref{prop:access-recoding} and the finite instruments
illustrate in their different settings. A memory label alone does not
alter a Cut, and the finite causal reader has not thereby realized
the exact chiral boundary.

\subsection{From finite use to a maintained boundary}

The finite stress-flow compiler and the ideal seed theorem provide
two parts of a possible interface: implementable tasks and a standard
limiting structure. Their joint maintenance remains the following
conjecture.

\begin{conjecture}[Maintained modular boundaries]\label{conj:maintained}
In a suitable local carrier and stationary reference state, a physical
record-and-control process can maintain a hierarchy of finite task
permissions whose modular covariance and nesting errors tend to zero as the
charged control, memory and preparation resources increase, and whose ideal
limit is a standard seed-generated Cut. Changing the retained control history
can change the maintained boundary and hence subsequent response and access.
\end{conjecture}
The single-orbit theorem supplies a nonempty ideal target, and
Corollary~\ref{cor:physical-finite-stress} supplies finite words on
bounded-energy inputs with retained references. A realization of
Conjecture~\ref{conj:maintained} must form and sustain the primitive
permissions with the same body's records, preparation and control
resources. It fails for a proposed process if covariance errors remain
bounded away from zero, if no admissible stationary limit is standard,
or if altered histories only relabel unchanged instruments. These
conditions concern access formation, not a consciousness score.

\subsection{A reflective bridge and its competing explanations}

A stronger proposal relates this maintained causal organization to a
subject's reflective access. Consciousness~III calls it renewable causal
self-acquaintance: acquired distinctions among changes in the world,
the interface and the reading of retained records must remain usable
for correction. \emph{The World as Imprint} adds historically formed
organization shared across tasks, whose self-evaluation regulates
inquiry and maintenance \cite{PaperCIII,WorldImprint}. These are two
aspects of a common candidate bridge, not independent confirmations.

The proposal has an explanatory direction. What presents as a stable
world for this process depends on its usable boundary; what it can
attribute to its own condition depends on a maintained model of that
access. A reflective perspective would consist in the continuing
organization of this participation, rather than a separate internal
spectator reading a complete copy of it. Finite records and
single-copy restrictions can limit what it describes about itself.
Modular time, the update clock of a self-model and experienced time
remain different quantities until a further relation is established.

A display-only account can reproduce routine confidence without making
those distinctions necessary for later correction. An ordinary causal
controller can reproduce the full functional organization. A theory
of experience that does not require reflective self-access may accept
the same functional results while rejecting a claim about phenomenal
presence. The strongest objection is therefore substantive: an
explanation of organized access can leave open why anything is
experienced \cite{Chalmers}. The mathematical correspondence constrains
the physical organization available to an explanation, but does not
eliminate that question.

The limits of an experimental bridge can be stated exactly. For a
common budget and protocol class \(\mathfrak P\), let
\(P_0^\pi,P_1^\pi\) be the complete record laws in two proposed
conditions. Protocols use acquired records and condition-independent
private randomness. Include acquired outcomes, operation labels, accessible
initial information and retained correlated tapes. Put
\[
 \delta_{\mathfrak P}=\sup_{\pi\in\mathfrak P}
                   \|P_0^\pi-P_1^\pi\|_{\rm TV}.
\]
The following is the interventional boundary proved in
Consciousness~III, applied to this comparison \cite{PaperCIII}.

\begin{corollary}[A limit on perspective tests]
\label{cor:perspective-tests}
Any equal-prior binary decision in this class has error at least
\((1-\delta_{\mathfrak P})/2\); any recorded statistic in \([0,1]\)
has expected values differing by at most \(\delta_{\mathfrak P}\).
In particular, complete interventional record equivalence leaves
no discriminator in this class, including through confidence or
reflective reports.
\end{corollary}
\begin{proof}
For each protocol, equal-prior testing error is bounded below by
one half minus one half of its total variation distance. That
distance also bounds differences of bounded statistics.
Taking the supremum over the same permitted class gives both
claims.
\end{proof}

This separates two empirical questions. A routine-matched scalar
calibrator can be distinguished from a controller retaining causal
access information by the correction task above. A controller with
the same complete interventional laws cannot be distinguished by
calling one of its internal variables first-person. Any proposed
experiential contrast needs an independently justified marker and
an intervention class in which its distribution actually differs.

\subsection{Discriminating conditions}

A useful experiment must keep first-order content performance,
routine confidence and the read/actuation budget matched before the
diagnostic event. It should then manipulate the world, the sensing
interface and retained-record reading separately, score the inferred
cause and actual correction, and compare acquired own-process
information with replay from independent trials. Correcting a report
is not the same task as restoring its physical cause.

The stronger shared-organization hypothesis requires transfer beyond
the calibration routine. A retained access distinction learned in one
task should alter relevant inquiry or correction in an independently
tested task using that same interface before new calibration is
supplied. Compare selective interruption of the record-to-operation
path with a sham intervention and a nonspecific processing-load
control. Compare erasure of raw logs after preserving the sufficient
causal state with erasure of that state while leaving an autobiography
display available. Restoring the causal state and update program in
replacement hardware supplies a control against an unjustified
privilege of the original memory medium.

Cause-specific action and any independently validated reflective
marker must be scored separately. Before applying such a marker,
validate its relation to reflective access while controlling for
first-order performance, motor capacity and reward sensitivity.
Specify a minimally detectable change and an equivalence margin in
advance. If usable cause-specific return is abolished while the
validated reflective marker remains equivalent to sham within that
margin, the proposed dependence of reflective access on that return
is refuted. A generic causal controller remains a complete
competitor for functional success; independent calibrators remain
a competitor if they explain all cross-task transfer.

The experiment tests whether a particular organization helps sustain
a perspective, not whether every system with a correctable boundary
has experience. If all permitted laws and markers are matched,
Corollary~\ref{cor:perspective-tests} applies. The bridge then needs
new evidence rather than a more expansive name for the same record.

Null-plane copies have separate transverse source and coupling
requirements \cite{CTT,CF}; Consciousness~II and HAFF-D address
those physical interfaces.

\section{Conclusions and open problems}\label{sec:open}

The standard Cut correspondence makes a boundary and its geometry
mutually constrained. The local stress-seed construction supplies
one such boundary from a specified carrier and modular evolution.
Its displacement relative to a fixed excitation changes recoverable
information even when the vacuum Cut triples are isomorphic.
Perspective therefore involves the source--boundary relation, not
only the name or isomorphism class of its reconstructed net.

For a finite process, retained history matters when it changes the
permissions and corrections through which that relation is used.
The self-description results make this claim testable and rule out
a confidence display as a sufficient account in specified tasks.
They also impose an exact observational boundary on stronger
interpretations. The remaining work is to realize maintained
boundaries in one physical history and to test their proposed role
in reflective organization. The main unresolved questions are:


\begin{enumerate}[leftmargin=2em,itemsep=2pt]
\item \textbf{Physical formation.} Derive the carrier, stationary reference and
availability of the primitive local current controls from a common
process, and test
Conjecture~\ref{conj:maintained}. The seed-generated standard Cut is
constructed in the specified model; primitive modular dynamics and a finite
maintenance protocol are not.
\item \textbf{General seeds and inequivalent Cuts.}
Conjecture~\ref{conj:T1} outside the positive log-bump cyclicity theorem;
whether a seed-generated standard triple differs from all ordinary translated
half-line triples. Modular-only generation is proved for that bump class,
not for every nonzero stress seed.
\item \textbf{Finite resources.} Determine explicit gate costs and convergence
rates for Corollary~\ref{cor:physical-finite-stress}. It already supplies
uniform approximability on bounded-energy inputs with arbitrary retained
references; it does not give an optimal budget or a maintained Cut limit.
\item \textbf{Transverse families and gravity.} Coupling between null-plane copies,
area and focusing; the selection of the algebra family and the physical
source enter here. An invisible source is not thereby absent; changing modular
data and changing actual stress are different tests. For ball
families in higher dimensions Consciousness~II identifies the variation of area
minus $(d-1)/R$ times volume as minus the weighted linearized Einstein charge.
On a conformal net in Minkowski space, its supplied
$\mathrm{MC}_\eta$ consistency law with the CHM bulk modular energy,
after the declared boundary-charge subtraction and under that paper's
source and functional assumptions, is equivalent to the linearized
Einstein equation~\cite{PaperCII};
null-plane couplings remain open.
\item \textbf{Selection.} Why this noise group; how unrecorded transformations are
aligned or reduced and at what cost; the imprint of the discarded group in the
sectors of the selected algebra (for SU(2) level one: the subcategory of Virasoro
sectors of lowest weight \(j^2\) contained in the vacuum representation).
\item \textbf{Unbounded generators.} Which non-multiplicative semigroups are compatible
with a strict scaling covariance \(E_t\sigma_s=\sigma_sE_{e^{2\pi s}t}\) is unclassified: it
is known that such a semigroup preserves the vacuum, that it is trivial if it commutes
with \(\sigma\), and that it is zero if its generator is bounded.
\end{enumerate}


\appendix

\section{Proofs for single-orbit generation}\label{app:singleorbit-proofs}
\subsection{Proof of Lemma~\ref{lem:flowclosure}}
\begin{proof}
These statements can be checked in any smooth coordinate chart covering the
compact supports. Finite-energy vectors form a common core on which
\(T(v)+T(w)=T(v+w)\); essential self-adjointness and the Trotter formula give
linear closure. Fixed-compact \(C^\infty\) convergence implies convergence of
the circle fields in the Fourier norm
\(\sum_n|\hat v_n|(1+|n|^{3/2})\). Carpi--Weiner's strong-resolvent result then
gives strong convergence of their exponential unitaries
\cite[Thm.~4.4, Prop.~4.5]{CW}.

For the bracket, put \(\varphi_t=\mathrm{Exp}(tv)\). Up to scalar phases,
\(V(\varphi_t)=e^{itT(v)}\in N\). Geometric covariance gives
\[
 V(\varphi_t)e^{iuT(w)}V(\varphi_t)^*
   =e^{iu\lambda(t,w)}e^{iuT((\varphi_t)_*w)},\qquad \lambda(t,w)\in\mathbb R .
\]
The Schwarzian term contributes only the displayed scalar; projective phases
do not change membership in a unital algebra. Thus
\((\varphi_t)_*w\) has all its flows in \(N\). Linear closure applies to
\((w-(\varphi_t)_*w)/t\), and the smooth pushforward formula gives
\[
 \left.\frac{d}{dt}(\varphi_t)_*w\right|_{t=0}
       =v'w-vw'=-[v,w].
\]
For sufficiently small \(|t|\) all supports stay in one fixed compact set.
The difference quotient converges in every \(C^k\) norm, so the preceding
strong-resolvent closure proves the bracket claim. No commutator of unbounded
operators on an unspecified domain is used.
\end{proof}

\subsection{Proof of Theorem~\ref{thm:singleorbit}}
\begin{proof}
In \(u=\log x\), the vector field \(f(x)\partial_x\) is \(g(u)\partial_u\).
A dilation by \(\lambda\) pushes \(g\) to \(g(u-\log\lambda)\).
Past modular times therefore give exactly \(g_c(u)=g(u-c)\) with
\(c=-2\pi s\ge0\). Injectivity of \(\arctan\) gives all stress flows of
\(g_c\) in \(N=N^-_{O_f}\), as in Appendix~\ref{app:hull}, Step~1.
All these fields have support to the right of \(\alpha\), proving the upper
inclusion in \eqref{eq:singleorbit}. For \(r\le0\),
\(\sigma_r(N)\subset N\); equivalently, already obtained fields can be shifted
further right by any \(c\ge0\).

Let \(d=\beta-\alpha\) and choose \(\varepsilon>0\) smaller than both endpoint
neighborhoods. Lemma~\ref{lem:flowclosure} supplies
\[
 k_\varepsilon=[g_0,g_{d-\varepsilon}],\qquad
 \operatorname{supp}k_\varepsilon\subset[\beta-\varepsilon,\beta].
\]
On the overlap, the first bump is nonincreasing and the second is
nondecreasing, both positive. Hence \(k_\varepsilon\ge0\).
Strict monotonicity need not mean a pointwise strict derivative. It does,
however, give a positive integral: the nonconstant increasing shifted bump
has a subinterval where its derivative is positive, and the other bump is
positive there. Thus
\[
 m_\varepsilon=\int k_\varepsilon(u)\,du>0,\qquad
 K_\varepsilon=k_\varepsilon/m_\varepsilon
\]
is a unit-mass positive kernel whose stress flows and all right translates'
flows belong to \(N\).

For \(h\in C_c^\infty((\beta,\infty),\mathbb R)\), set
\begin{equation}\label{eq:positivekernel}
 h_\varepsilon(u)=\int_{c\ge0}h(c+\beta)K_\varepsilon(u-c)\,dc
                =\int K_\varepsilon(w)h(u-w+\beta)\,dw .
\end{equation}
The second equality extends the integral without adding negative shifts:
the support of \(h(c+\beta)\) already has \(c>0\).
If its distance from \(\beta\) is \(\eta>0\), then for
\(\varepsilon<\eta/2\) the supports of \(h_\varepsilon\) lie in a fixed
compact subset of \((\beta,\infty)\). Since
\(w\in[\beta-\varepsilon,\beta]\) and the mass is one, for every \(j\ge0\),
\[
 \|h_\varepsilon^{(j)}-h^{(j)}\|_\infty
 \le\varepsilon\|h^{(j+1)}\|_\infty .
\]
The integral is a \(C^\infty\) limit of finite real Riemann sums of permitted
right-translated kernels, all with fixed compact support. Lemma~\ref{lem:flowclosure}
first gives the flows of \(h_\varepsilon\), then those of \(h\), in \(N\).
Returning through \(x=e^u\), these are all compact stress fields on
\((b,\infty)\). They generate \(\mathrm{Vir}_c(b,\infty)\), giving the lower
inclusion.

Reeh--Schlieder makes \(\Omega\) cyclic for the lower half-line algebra, hence
for \(N\); it is separating for \(N\subset M\). Locality gives
\(\mathrm{Vir}_c(0,a)\subset M\cap N'\), so the relative commutant is also
vacuum-cyclic. The upper half-line is proper: otherwise the nontrivial
\(\mathrm{Vir}_c(0,a)\) would lie in the center of the factor \(M\).
The past-orbit definition supplies half-sided invariance.
Theorems~\ref{thm:AZ} and~\ref{thm:GLW} reconstruct translations and the net.
\end{proof}

\section{Proof of Theorem~\ref{thm:hull}}\label{app:hull}

Write \(B=\mathrm{Vir}_c(0,\infty)\), \(\mathrm{Hull}=\mathrm{Hull}_{\sigma,T}(O_f)\), and for a
von Neumann algebra \(N\subset M\)
\(\mathcal V_N=\{h\in C_c^\infty((0,\infty),\mathbb R):e^{iuT(h)}\in N\ \forall u\in\mathbb R\}\).

\paragraph{Step 0: \(\mathrm{Hull}\subset B\).} \(B\) is admissible: \(\sigma_s=\mathrm{Ad}\,V(\delta_{e^{-2\pi s}})\) preserves
the supports \((0,\infty)\) and the generators \(T(h)\) go to \(T(\delta_*h)\); \(T_a\), \(a\ge0\),
moves supports into \((a,\infty)\subset(0,\infty)\); and \(O_f\in B\).

\paragraph{Step 1: spectral projections.} \(O_f=\psi(T(f))\) with the injective Borel function
\(\psi=\tfrac2\pi\arctan\). Hence \(\psi^{-1}\) is Borel, every bounded Borel function of
\(T(f)\) is a bounded Borel function of \(O_f\), and so \(e^{iuT(f)}\in W^*(O_f)\subset N\)
for every admissible \(N\ni O_f\). In particular \(f\in\mathcal V_N\); \(\mathcal V_N\) is
closed under multiplication by real scalars since \(u\) ranges over \(\mathbb R\).

\paragraph{Step 2: affine images.} For \(\varphi(x)=\lambda x+b\), \(\lambda>0\), \(b\ge0\), the Schwarzian
vanishes, so \(V(\varphi)T(h)V(\varphi)^*=T(\varphi_*h)\) with
\((\varphi_*h)(y)=\lambda h((y-b)/\lambda)\). If \(h\in\mathcal V_N\) then
\(e^{iuT(\varphi_*h)}=T_b\big(\sigma_s(e^{iuT(h)})\big)\) with \(e^{-2\pi s}=\lambda\)
belongs to \(N\), because \(N\) is stable under \(\sigma\) and under \(T_b\). So
\(\varphi_*h\in\mathcal V_N\).

\paragraph{Step 3: linear and limit closure.} If \(h_1,h_2\in\mathcal V_N\) then
\(h_1+h_2\in\mathcal V_N\): \(T(h_1)+T(h_2)\) equals \(T(h_1+h_2)\) on the finite-energy
vectors, a common core on which \(T(h_1+h_2)\) is essentially self-adjoint
\cite[Thm.~4.4]{CW}, so the Trotter product formula
\(e^{iu(A+B)}=\text{s-}\lim(e^{iuA/n}e^{iuB/n})^n\) gives
\(e^{iuT(h_1+h_2)}\in N\), \(N\) being strongly closed. If \(h_n\in\mathcal V_N\) have
supports in a fixed compact subset of \((0,\infty)\) and \(h_n\to h\) in \(C^\infty\)
(hence in \(\|\cdot\|_{3/2}\)), then \(e^{iuT(h_n)}\to e^{iuT(h)}\) strongly
\cite[Prop.~4.5]{CW}, and \(h\in\mathcal V_N\).

\paragraph{Step 4: the case \(m=\int f\ne0\).} Let \(h\in C_c^\infty((0,\infty),\mathbb R)\)
and put \(f_{\varepsilon,b}(y)=\varepsilon f((y-b)/\varepsilon)\), \(k_{\varepsilon,b}=f_{\varepsilon,b}/(\varepsilon^2m)\), so
\(\int k_{\varepsilon,b}=1\) and \(k_{\varepsilon,b}\in\mathcal V_N\) by Steps 1--3. Then
\[
h_\varepsilon(y):=\int_{b\ge0} h(b)k_{\varepsilon,b}(y)\,db=\frac1m\int f(w)\,h(y-\varepsilon w)\,dw\longrightarrow h(y)
\]
in \(C^\infty\), with supports in a fixed compact subset of \((0,\infty)\) for small
\(\varepsilon\). Only \(b\ge0\) is used: \(h\) has compact positive support, and
\(h(y-\varepsilon w)=0\) whenever \(b=y-\varepsilon w<0\).
Riemann sums of the integral are finite real combinations of the
\(k_{\varepsilon,b_i}\), which lie in \(\mathcal V_N\) and converge to \(h_\varepsilon\) in
\(C^\infty\) with fixed compact support, so \(h_\varepsilon\in\mathcal V_N\) by Step 3, and
then \(h\in\mathcal V_N\). Hence \(\mathcal V_N=C_c^\infty((0,\infty),\mathbb R)\).

\paragraph{Step 5: the case \(\int f=0\).} Let \(\varphi_s=\mathrm{Exp}(sf)\). By Step 1,
\(e^{isT(f)}\in N\), and \(e^{isT(f)}=V(\varphi_s)\) up to a phase, so conjugation by it
maps \(\mathcal V_N\) to itself via \(h\mapsto(\varphi_s)_*h\). (For a general diffeomorphism the conjugation
sends \(T(h)\) to \(T(\varphi_*h)\) plus a real multiple of the identity, the Schwarzian term; the
multiple exponentiates to a phase, which does not affect membership in \(N\) since \(N\ni1\).) With
\(f_\delta(y)=f(y-\delta)\), \(\delta\ge0\), we get \((\varphi_s)_*f_\delta\in\mathcal V_N\).
Since \(\tfrac d{ds}(\varphi_s)_*g=-[f,g]\) with \([X,Y]=XY'-X'Y\) and
\(\int[f,f_\delta]=2\int ff_\delta'\),
\[
\frac d{ds}\Big|_{0}\int(\varphi_s)_*f_\delta\,dy=-2\int f(y)f'(y-\delta)\,dy=2C'(\delta),
\qquad C(\delta)=\int f(y)f(y-\delta)\,dy .
\]
\(C\) is even, smooth, compactly supported and \(C(0)=\|f\|^2>0\), so \(C'(\delta_0)\ne0\)
for some \(\delta_0>0\). For small \(s\ne0\) the field \(g=(\varphi_s)_*f_{\delta_0}\) has
\(\int g\ne0\), compact support in \((0,\infty)\) and \(g\in\mathcal V_N\); apply Step 4 to \(g\).

\paragraph{Step 6.} All \(e^{iuT(h)}\), \(h\in C_c^\infty((0,\infty),\mathbb R)\), lie in \(N\), so
\(N\supset\{T(h)\}''=B\). With Step 0, \(\mathrm{Hull}=B\). \(\square\)

\section{The Virasoro vacuum module in Fock space}\label{app:fock}

In the Fock space of the \(U(1)\) current with \([J_m,J_n]=m\delta_{m+n,0}\),
\(J_n\Omega=0\) for \(n\ge0\), take \(L_n=\tfrac12\sum_k{:}J_{n-k}J_k{:}\). We built the descendants \(L_{-n_1}\cdots L_{-n_k}\Omega\) (\(n_i\ge2\)) at each level with exact
rational arithmetic and computed their rank, and compared with the number of
partitions into an even number of parts (the \(\mathbb Z_2\)-even Fock space).
\begin{center}
\begin{tabular}{l|cccccccccc}
level \(n\) & 0&1&2&3&4&5&6&7&8&9\\\hline
\(\dim\) Virasoro vacuum module & 1&0&1&1&2&2&4&4&7&8\\
\(p(n)-p(n-1)\) & 1&0&1&1&2&2&4&4&7&8\\
\(\dim\) even Fock space & 1&0&1&1&3&3&6&7&12&14
\end{tabular}
\end{center}
The first level where the even space is strictly larger is 4. At level 4,
\(\langle L_{-4}\Omega,L_{-4}\Omega\rangle=5=\tfrac c{12}(4^3-4)\) with \(c=1\), the Gram
matrix on \(\{L_{-4}\Omega,L_{-2}^2\Omega\}\) is \(\left(\begin{smallmatrix}5&3\\3&9/2\end{smallmatrix}\right)\) with
determinant \(27/2\), and
\(u=\tfrac12J_{-1}^4\Omega-J_{-3}J_{-1}\Omega+\tfrac34J_{-2}^2\Omega\) satisfies
\(L_1u=L_2u=L_3u=L_4u=0\), \(\langle u,L_{-4}\Omega\rangle=\langle u,L_{-2}^2\Omega\rangle=0\),
\(\|u\|^2=\tfrac14\cdot24+3+\tfrac9{16}\cdot8=\tfrac{27}2\); the three vectors have rank
3. The generating function of the even dimensions agrees through order 23 with
\(\sum_{p\ge0}(t^{4p^2}-t^{(2p+1)^2})\prod_n(1-t^n)^{-1}\), and the \(SU(2)_1\) vacuum character
identity \(\sum_mz^mt^{m^2}P(t)=\sum_j\chi_j(z)(t^{j^2}-t^{(j+1)^2})P(t)\), \(P(t)=\prod(1-t^n)^{-1}\),
holds coefficient by coefficient for \(|m|\le5\) through order 23. These finite checks are consistent
with the decomposition of the even Fock space as \(\bigoplus_pL(1,4p^2)\) and with the character
identity behind Theorem~\ref{thm:Rehren}, but a check through finitely many levels does not prove an
infinite decomposition; only the level-4 computation is used in Proposition~\ref{prop:U1}, and the
identity of characters used for Theorem~\ref{thm:Rehren} is Rehren's.

\section{Numerical checks}\label{app:num}

\paragraph{Cut displacement.} For \(f(x)=e^{4-1/(z(1-z))}(1+0.7(z-\tfrac12))\), \(z=x-2\), supported in
\((2,3)\), \(c=1\), \(\rho=\mathrm{Exp}(\theta f)\), the integrals of Theorem~\ref{thm:displacement}
were evaluated by quadrature after integrating the flow and its first three
variational derivatives. For \(\theta=0.05\):
\begin{center}
\begin{tabular}{r|cccc}
\(a\) & 0 & 1 & 1.5 & 1.9\\\hline
\(D_a/D_0\) & 1.000 & 0.604 & 0.406 & 0.2476\\
\(D_0-D_a\) & 0 & 4.5641e-2 & 6.8461e-2 & 8.6718e-2\\
\(2\pi aE\) & 0 & 4.5641e-2 & 6.8461e-2 & 8.6718e-2
\end{tabular}
\end{center}
and both \(-\partial_aD_a\) and \(2\pi\int\langle T\rangle dx\) are
\(4.5640958\times10^{-2}\). At 2001 and 4001 spatial points the displayed
values agree to the shown precision. This is a quadrature diagnostic, not an
interval-certified error bound.
For \(\theta=0.2\), \(D_a/D_0=1,0.593,0.390,0.227\) at the same \(a\).

\paragraph{Single-orbit closure.} The analytic proof uses the allowed
kernel integral \eqref{eq:positivekernel}, without negative recentering.
For \(g(u)=e^{4-1/(u(1-u))}\) on \((0,1)\), numerical quadrature checks
the kernel's nonnegative sign and positive mass, and convergence in the first
three derivative norms for targets supported strictly to the right of \(1\).
A separate geometric-flow check evaluates the pushforward difference quotient
of Lemma~\ref{lem:flowclosure}. These diagnostics check signs and a finite
selection of derivatives; the proof of \(C^\infty\) and operator closure is
the analytic argument, not the grid. Small kernel mass is also recorded,
since normalization can be extremely costly and cannot certify finite-resource
generation.

\paragraph{Finite-dimensional checks.} The identities of Theorems~\ref{thm:four} and
\ref{thm:multi} were checked in the spin representation \([1,1,2]\) of \(SO(3)\) on
\(M_{11}\): the fixed algebra is 5-dimensional for the two-axis noise, the axis noises
do not commute (\(\|[Q_z,Q_x]\|_F=1.90\)), the dissipative identity of
Theorem~\ref{thm:multi} holds for a fixed complex random probe
(both sides \(363.2646789268\)), and for \(L=Q_x+Q_z-2\,\mathrm{id}\),
\(\|e^{tL}-Q_{SO(3)}\|\) falls from \(0.60653\) at \(t=1\) to
\(9.4\times10^{-14}\) at \(t=60\). The two specified irrational rotations
also have a five-dimensional fixed algebra in this representation. These
are finite-dimensional analogues, not statements about the net.

\section{Boundaries of shell and scale rules}\label{app:july}

A shell containing a full norm ball generates the ambient algebra and
cannot select a proper Cut; a unit-sphere version has the same defect.
Modular covariance must respect the antilinearity identity
$J\Delta^{is}J=\Delta^{is}$. The vacuum modular generator has zero
variance, while a truncated one-sided variance can diverge and have a
regularization-dependent finite part. Approximate innerness does not
imply innerness of an automorphism. Faithfulness gives separation in
the stated representation, not cyclicity.
These boundaries are separate from the local stress-seed proof.
The complete earlier-route analysis is retained in the supplementary
source record.

\begin{thebibliography}{99}
\bibitem{QRC} S.~Liu, \emph{Representing Control in the World: Retained History, Physical Sources and Renewal in Q-RAIF}, companion manuscript (2026).
\raggedright

\bibitem{AZ} H.~Araki and L.~Zsid\'o, Extension of the structure theorem of Borchers
and its application to half-sided modular inclusions, \emph{Rev.\ Math.\ Phys.}
\textbf{17} (2005) 491--543; arXiv:math/0412061.
\bibitem{BGL} R.~Brunetti, D.~Guido and R.~Longo, Modular structure and duality in
conformal quantum field theory, \emph{Commun.\ Math.\ Phys.} \textbf{156} (1993) 201--219.
\bibitem{BKK} C.~B\'eny, A.~Kempf and D.~W.~Kribs, Generalization of quantum error
correction via the Heisenberg picture, \emph{Phys.\ Rev.\ Lett.} \textbf{98} (2007)
100502; arXiv:quant-ph/0608071.
\bibitem{BW} J.~J.~Bisognano and E.~H.~Wichmann, On the duality condition for a Hermitian
scalar field, \emph{J.\ Math.\ Phys.} \textbf{16} (1975) 985.
\bibitem{BSM} D.~Buchholz and H.~Schulz-Mirbach, Haag duality in conformal quantum field theory, \emph{Rev.\ Math.\ Phys.} \textbf{2} (1990) 105--125.
\bibitem{Carpi} S.~Carpi, On the representation theory of Virasoro nets,
\emph{Commun.\ Math.\ Phys.} \textbf{244} (2004) 261--284; arXiv:math/0306425.
\bibitem{Chalmers} D.~J.~Chalmers, Facing up to the problem of consciousness, \emph{J.\ Consciousness Stud.} \textbf{2} (1995) 200--219.
\bibitem{CR} A.~Connes and C.~Rovelli, Von Neumann algebra automorphisms and time-thermodynamics relation in generally covariant quantum theories,
\emph{Class.\ Quantum Grav.} \textbf{11} (1994) 2899--2917.
\bibitem{CW} S.~Carpi and M.~Weiner, On the uniqueness of diffeomorphism symmetry in
conformal field theory, \emph{Commun.\ Math.\ Phys.} \textbf{258} (2005) 203--221; arXiv:math/0407190.
\bibitem{CF} Y.~Ceyhan and T.~Faulkner, Recovering the QNEC from the ANEC,
\emph{Commun.\ Math.\ Phys.} \textbf{377} (2020) 999--1045; arXiv:1812.04683.
\bibitem{CTT} H.~Casini, E.~Teste and G.~Torroba, Modular Hamiltonians on the null
plane and the Markov property of the vacuum state, \emph{J.\ Phys.\ A} \textbf{50}
(2017) 364001; arXiv:1703.10656.
\bibitem{GW} R.~Goodman and N.~R.~Wallach, Projective unitary positive-energy representations of Diff(\(S^1\)), \emph{J.\ Funct.\ Anal.} \textbf{63} (1985) 299--321.
\bibitem{GLW} D.~Guido, R.~Longo and H.-W.~Wiesbrock, Extensions of conformal nets and
superselection structures, \emph{Commun.\ Math.\ Phys.} \textbf{192} (1998) 217--244;
arXiv:hep-th/9703129.
\bibitem{Hollands} S.~Hollands, Relative entropy for coherent states in chiral CFT,
arXiv:1903.07508.
\bibitem{KR} V.~G.~Kac and A.~K.~Raina, \emph{Highest Weight Representations of Infinite Dimensional Lie Algebras}, World Scientific, 1987.
\bibitem{Rehren} K.-H.~Rehren, News from the Virasoro algebra (published as: A new view of
the Virasoro algebra), \emph{Lett.\ Math.\ Phys.} \textbf{30} (1994) 125--130;
arXiv:hep-th/9308144.
\bibitem{Segal} G.~Segal, Unitary representations of some infinite dimensional groups, \emph{Commun.\ Math.\ Phys.} \textbf{80} (1981) 301--342.
\bibitem{PaperCII} S.~Liu, \emph{Higher-Dimensional Consistency of Geometry and Modular Flow: Physical Source Identification and a Second-Order Extension} (Consciousness~II), 2026;
July archive DOI: 10.5281/zenodo.19607543.
\bibitem{FoundationI} S.~Liu, \emph{Foundation I: From Source to Manifestation.
Accessible Algebras from a Source, Its Noise and a Readable Seed},
2026; July archive DOI: 10.5281/zenodo.19627786.
\bibitem{PaperA} S.~Liu, \emph{Physical Coarse-Graining and the Problem of Emergent Geometry:
The Holographic Alaya--Field Framework} (HAFF-A), 2026.
\bibitem{PaperB} S.~Liu, \emph{Coupling, Records, and Stable Relations:
Formation, Protected Content and the Lifetime of a Geometry} (HAFF-B),
2026.
\bibitem{PaperE} S.~Liu, \emph{Measurement as Accessibility:
Measurement Permissions, Stable Records and Public Readout} (HAFF-E),
2026.
\bibitem{Takesaki} M.~Takesaki, Conditional expectations in von Neumann algebras, \emph{J.\ Funct.\ Anal.} \textbf{9} (1972) 306--321.
\bibitem{Xu} F.~Xu, Strong additivity and conformal nets, \emph{Pacific J.\ Math.}
\textbf{221} (2005) 167--199; arXiv:math/0303266.

\bibitem{PaperCIII} S.~Liu, \emph{Finite Self-Models and the Explanatory Gap:
What a Continuing Observer Can Know About Itself}
(Consciousness~III), 2026.
\bibitem{WorldImprint} S.~Liu, \emph{The World as Imprint:
From Trace-Carrying Structure to Self-Reading Mind}, 2026.

\end{thebibliography}

\end{document}
